% Editable CNN comparison: which input pixels one output pixel depends on,
% for a positionwise MLP and for a fully-connected layer on the vectorized image.
% Redrawn by Claude (2026-10-15) after Simone's review ("too basic").
\begingroup
% #1,#2: lower-left corner; #3: 1 = only the centre pixel is used,
% 2 = every pixel is used, 0 = output grid (centre pixel highlighted).
\newcommand{\cnnPixelGrid}[3]{%
    \begin{scope}[shift={(#1,#2)}]
        % two offset layers suggest the channel axis
        \foreach \k in {2,1}{
            \fill[drawgray_l] (0.12*\k,0.12*\k) rectangle ++(1.65,1.65);
            \draw[drawgray!45,line width=0.4pt] (0.12*\k,0.12*\k) rectangle ++(1.65,1.65);
        }
        \foreach \row in {0,1,2}{
            \foreach \col in {0,1,2}{
                \pgfmathtruncatemacro{\centre}{(\row==1)&&(\col==1)}
                \ifnum#3=2
                    \fill[drawblue_l] (0.55*\col,0.55*\row) rectangle ++(0.55,0.55);
                \else
                    \fill[white] (0.55*\col,0.55*\row) rectangle ++(0.55,0.55);
                    \ifnum\centre=1
                        \ifnum#3=1
                            \fill[drawblue_l] (0.55*\col,0.55*\row) rectangle ++(0.55,0.55);
                        \else
                            \fill[drawgreen_l] (0.55*\col,0.55*\row) rectangle ++(0.55,0.55);
                        \fi
                    \fi
                \fi
                \draw[drawgray!70,line width=0.4pt] (0.55*\col,0.55*\row) rectangle ++(0.55,0.55);
            }
        }
        \ifnum#3=0
            \draw[drawgreen,line width=1.1pt] (0.55,0.55) rectangle ++(0.55,0.55);
        \else
            \ifnum#3=1
                \draw[drawblue,line width=1.1pt] (0.55,0.55) rectangle ++(0.55,0.55);
            \else
                \draw[drawblue,line width=1.1pt] (0,0) rectangle ++(1.65,1.65);
            \fi
        \fi
    \end{scope}
}

\begin{tikzpicture}[x=1cm,y=1cm,
    block/.style={draw=drawgreen,fill=drawgreen_l!40!white,rounded corners=2pt,
        minimum width=1.5cm,minimum height=0.8cm,inner sep=3pt,font=\normalsize},
    wire/.style={drawblue,line width=0.7pt},
    edge/.style={-{Stealth[length=2mm]},line width=0.8pt,draw=drawgray},
    heading/.style={font=\normalsize\bfseries,anchor=west,inner sep=0pt},
    note/.style={font=\small,anchor=west,align=left,inner sep=0pt}]

% ---- Positionwise MLP: one input pixel per output pixel
\node[heading] at (0,5.05) {Positionwise MLP};
\cnnPixelGrid{0}{2.9}{1}
\cnnPixelGrid{6.2}{2.9}{0}
\node[block] (pw) at (4.0,3.725) {$f_{\boldsymbol{\theta}}$};
\fill[drawblue] (0.825,3.725) circle (1.6pt);
\draw[wire] (0.825,3.725) -- (1.65,3.725) -- (pw.west);
\draw[edge] (pw.east) -- (6.75,3.725);
\node[font=\footnotesize] at (0.825,2.6) {$X\sim(H,W,C)$};
\node[font=\footnotesize] at (7.025,2.6) {$Y\sim(H,W,C')$};
\node[note] at (8.6,4.05) {$\idx{Y}{i,j,:}$ sees only $\idx{X}{i,j,:}$};
\node[note,text=drawgray] at (8.6,3.45) {weights do not depend on $H,W$};

\draw[drawgray!35,line width=0.5pt] (0,2.15)--(13.4,2.15);

% ---- Vectorized image: every input pixel for every output pixel
\node[heading] at (0,1.75) {Vectorized image};
\cnnPixelGrid{0}{-0.4}{2}
\cnnPixelGrid{6.2}{-0.4}{0}
\node[block] (fc) at (4.0,0.425) {$\vc{W}$};
% one wire per pixel, fanning out from the right edge of the grid
\foreach \row in {0,1,2}{
    \foreach \col in {0,1,2}{
        \fill[drawblue] (0.275+0.55*\col,-0.125+0.55*\row) circle (1.6pt);
    }
}
\foreach \k in {0,...,8}{
    \draw[wire] (1.65,-0.32+0.18*\k) -- (fc.west);
}
\draw[edge] (fc.east) -- (6.75,0.425);
\node[note] at (8.6,0.75) {$\idx{Y}{i,j,:}$ sees every pixel of $X$};
\node[note,text=drawgray] at (8.6,0.15) {a layer has $HWC\cdot HWC'$ weights};
\end{tikzpicture}
\endgroup
